% sections/09_conclusion.tex

\section{Conclusion}
\label{sec:conclusion}

We presented \sysname, a publish-time static \emph{risk-prioritization} workflow for
surfacing capability-description mismatches in LLM agent skill marketplaces.
The workflow combines deterministic multi-file capability extraction (LSCPS,
14 families, \ResultFilesRead{} artifact files across \NskillsDMR{} skills)
with a lightweight disclosure-matching layer that computes Capability Recall
(CR), Description Precision (CP), and Capability Discrepancy Index (CDI).

Our evaluation on \NskillsDMR{} \clawhub{} skills finds that
\ResultICone{} of skills exhibit complete capability non-disclosure (IC-1,
CR\,=\,0) and \ResultICtwo{} exhibit partial non-disclosure (IC-2),
yielding corpus-mean CR\,=\,\ResultCR{} and CDI\,=\,\ResultCDI{}.
These findings are robust to regex bank variation: under 8 perturbation conditions
(strict/relaxed thresholds, 4 family ablations, 5-family taxonomy collapse),
inconsistency rates remain structurally elevated, with the Strict condition
still showing $>$47\% total inconsistency.
The combined SCS\,$\cup$\,DMR rule retrieves 94.1\%
(95\% Wilson CI:\,[80.9\%,\,98.4\%]) of \Nmalicious{} confirmed malicious skills
that the platform's own scanner missed entirely,
and outperforms four natural NLP baselines (lexical overlap, TF-IDF distance,
single-file static, general-purpose linters) on every ranking metric
(MRR: 0.019 vs. $\leq$0.002 for all baselines).
A marketplace security team reviewing the top-200 \sysname-ranked skills
recovers 17.6\% of confirmed T6 cases with a 16.3$\times$ lift over random,
compared to 0\% for the platform scanner under any budget.
Cross-platform validation on \NskillsSH{} independent skills suggests
that capability-description inconsistency exceeds 90\%
across both ecosystems without parameter retuning.

\paragraph{Limitations.}
LSCPS is a conservative lower bound: capabilities expressed via
non-standard idioms or obfuscated code are not captured.
Correspondingly, CR\,=\,\ResultCR{} should be interpreted as an
upper bound on actual disclosure---true gaps may be smaller.
The T6 ground truth (\Nmalicious{} skills) is limited by available
confirmed cases; recall estimates carry wide confidence intervals.
The \clawhub{} ecosystem, while the largest LLM skill marketplace
at the time of study, is smaller than mature package registries
(PyPI: $>$500K packages; npm: $>$2M packages).
We emphasize \emph{early-warning value}: documenting and hardening
the threat surface at 18K skills---when structural remediation
(per-skill capability manifests, YAML declaration enforcement)
is still tractable---is precisely when measurement papers are most actionable.
The npm event-stream and node-ipc precedents (\S\ref{sec:related})
demonstrate that supply-chain attacks scale faster than defenses;
our goal is to break that pattern before it repeats in LLM skill ecosystems.

\paragraph{Future Work.}
\textbf{AutoDesc and CapManifest} (§\ref{sec:autodesc}) represent the
primary constructive extension: an LLM-assisted description drafting tool
driven by LSCPS extraction results, paired with a machine-readable
CapManifest protocol that closes the IC gap from the developer side.
Evaluation goals include CR improvement rate ($\text{CR}_{\text{after}} \to 1.0$),
description acceptability in user studies, and false-capability rate.
\textbf{LLM semantic adjudication:} the five-model open-source ensemble
(prototype implemented) would reduce LSCPS false positives by verifying
that matched patterns are semantically functional.
\textbf{Expanded T6 ground truth} via coordinated platform disclosure
would improve statistical power on recall estimates.
\textbf{Real-time deployment}: integrating SCS+DMR as a publish-time
CI/CD gate would move \sysname from offline triage to active marketplace governance.
We view \sysname and AutoDesc as complementary layers---detection and remediation---
that together make the path of accurate disclosure easier than the path of omission.

\section*{Acknowledgments}

\textit{Anonymized for review.}

